% Provenance (never printed): see sources/outline_sources.md (Year 3:
% partition coefficients; batch and fractional melting; crystallisation
% models; reading a spider diagram quantitatively). Brief:
% sources/book5/BRIEFS.md, ch. 12. Curves from figdata/bachelor-3/
% melting_models.py. Exercise and problem values in the Book 5 scratch
% calculator calc/ch12.py. Aldren Island and its lavas are invented.
\chapter{Modélisation des éléments en traces}\label{ch:b3:trace-element-modelling}

Deux basaltes ont fait éruption sur la même île, à quelques kilomètres l'un de l'autre. L'un
contient plus d'une fois et demie autant de lanthane que l'autre, et presque quatre fois
moins d'ytterbium. Ont-ils fondu à partir de deux manteaux différents, ou à partir d'un même manteau
fondu à des taux différents, ou à des profondeurs différentes~? Les éléments majeurs
ne permettent guère de trancher~: ce sont deux basaltes. Les éléments en traces, qui suivent les minéraux
sans les modifier, conservent la trace de la quantité de liquide produite, de la manière dont il a quitté
sa source et quels minéraux sont restés en place. Ce chapitre écrit cette trace
sous forme d'équations et la relit.

\begin{recall}
Le coefficient de partage minéral--liquide $D$, le coefficient de partage global,
les éléments compatibles et incompatibles, les terres rares, la chondrite CI
et le manteau primitif, les diagrammes normalisés et l'anomalie en
europium, MORB et OIB figurent dans le volume de deuxième année, dans son chapitre sur les éléments dans la Terre.
La péridotite, la lherzolite, la fusion partielle et le taux de fusion $F$
figurent dans son chapitre sur la fusion du manteau~; la cristallisation à l'équilibre et
fractionnée, la loi de fractionnement de Rayleigh, le mélange et
l'assimilation--cristallisation fractionnée (AFC), dans son chapitre sur la
différenciation magmatique. Y a été démontrée la constance du rapport de deux éléments très
incompatibles pendant la cristallisation.
\end{recall}

\section{La comptabilité d'une roche en fusion}\label{sec:b3:trace-element-modelling:balance}

Une roche de masse 1 fond. À la fraction de fusion $F$, elle a produit une masse $F$ de
liquide et conservé une masse $1 - F$ de résidu solide. Un élément en traces présent à la
concentration $C_0$ dans la source se répartit entre les deux~; sa concentration
dans le liquide est $C_L$, dans le résidu $C_S$. Chaque modèle de ce chapitre est une
forme d'un même bilan, $C_0 = FC_L + (1 - F)C_S$ pour l'ensemble du liquide et du résidu,
et d'une règle qui relie $C_S$ à $C_L$~: au contact entre un cristal et
le liquide, $C_{\text{minéral}} = D\,C_L$. Les modèles diffèrent par ce qui touche
à quoi, et à quel moment.

Les coefficients de partage des minéraux du manteau varient largement. Dans une
compilation, le lanthane a $D = 0.058$ dans le clinopyroxène et $0.001$
dans le grenat, dans l'olivine, l'orthopyroxène et le spinelle~; l'ytterbium a $D = 0.58$
dans le clinopyroxène, $4.5$ dans le grenat et $0.16$ dans les trois autres. Le grenat retient les
terres rares lourdes et rejette les légères.
% ledger: budahn91:d-ree

\begin{definition}[Fusion non modale]\label{def:b3:trace-element-modelling:nonmodal}
Dans la \emph{fusion non modale}\index{fusion non modale}, les minéraux de la
source passent dans le liquide selon des proportions $p_i$ qui diffèrent de leurs fractions
massiques $x_i$ dans la source. Le liquide voit alors le coefficient de partage global
de la source, $D_0 = \sum_i x_iD_i$, au début de la fusion, puis le
coefficient pondéré par la fusion $P = \sum_i p_iD_i$ à mesure que celle-ci progresse.
\end{definition}

La fusion est rarement modale. Dans un exemple travaillé de lherzolite à spinelle, composée de
\qty{58}{\percent} d'olivine, \qty{27}{\percent} d'orthopyroxène,
\qty{12}{\percent} de clinopyroxène et \qty{3}{\percent} de spinelle, les minéraux
passent dans le liquide selon les proportions 20, 25, 45 et \qty{10}{\percent}~:
le clinopyroxène fond quatre fois plus vite que sa part, l'olivine trois fois
plus lentement.
% ledger: white13:ex73

\begin{proposition}[Coefficient de partage global du résidu]
\label{prop:b3:trace-element-modelling:residue}
Dans la \omterm{def:b3:trace-element-modelling:nonmodal}{fusion non modale} à $p_i$ constants, le résidu à la fraction de fusion $F$
a pour coefficient de partage global
\[
\bar D(F) = \frac{D_0 - PF}{1 - F},
\]
valable tant que chaque minéral subsiste, c'est-à-dire pour $F < x_i/p_i$ pour chaque $i$.
\end{proposition}

\begin{proof}
Par unité de masse de source, la masse du minéral $i$ restant à $F$ est $x_i - p_iF$,
et le résidu a pour masse $1 - F$. Son coefficient global est la moyenne pondérée par les
masses, $\sum_i (x_i - p_iF)D_i/(1 - F) = (D_0 - PF)/(1 - F)$. Un minéral est
épuisé lorsque $x_i - p_iF = 0$.
\end{proof}

Avec les coefficients minéraux ci-dessus, une lherzolite à spinelle de l'exemple
travaillé, et une lherzolite à grenat maintenue à \qty{5}{GPa} et
\qty{1400}{\celsius} en laboratoire, de composition \qty{54.1}{\percent} d'olivine,
\qty{9.8}{\percent} d'orthopyroxène, \qty{19.5}{\percent} de clinopyroxène et
\qty{16.7}{\percent} de grenat en masse, on obtient les coefficients globaux ci-dessous.
% ledger: white13:ex73, rfr25:gkr-5gpa, budahn91:d-ree
\begin{center}
\small
\begin{tabular}{lccc}
\hline
source & $D_{\mathrm{La}}$ & $D_{\mathrm{Sm}}$ & $D_{\mathrm{Yb}}$ \\
\hline
lherzolite à spinelle & 0.0078 & 0.062 & 0.21 \\
lherzolite à grenat & 0.012 & 0.107 & 0.97 \\
\hline
\end{tabular}
\end{center}
Le lanthane est très incompatible dans les deux. L'ytterbium est modérément incompatible
dans la lherzolite à spinelle et presque neutre dans la lherzolite à grenat. La
différence entre les deux sources tient presque entièrement aux terres rares
lourdes.

\section{Fusion à l'équilibre}\label{sec:b3:trace-element-modelling:batch}

\begin{theorem}[Fusion à l'équilibre]\label{thm:b3:trace-element-modelling:batch}
Dans la \emph{fusion à l'équilibre}\index{fusion à l'équilibre}, le liquide reste en équilibre
avec l'ensemble du résidu jusqu'à ce qu'il quitte la source d'un seul coup. Avec un coefficient
de partage global $D$ constant, le liquide et le résidu à la fraction de fusion $F$
vérifient
\[
\frac{C_L}{C_0} = \frac{1}{D + F(1 - D)}, \qquad
\frac{C_S}{C_0} = \frac{D}{D + F(1 - D)}.
\]
\end{theorem}

\begin{proof}
À l'équilibre, chaque cristal du résidu contient $D_iC_L$, si bien que le
résidu contient en moyenne $C_S = DC_L$. Le bilan $C_0 = FC_L + (1 - F)DC_L$ donne
$C_L$, et $C_S$ en découle.
\end{proof}

La formule est la loi de cristallisation à l'équilibre du volume de deuxième année lue
à rebours~: dans les deux cas, $F$ est la fraction de liquide dans un système clos. Ses deux
limites méritent d'être connues. Pour $F \to 0$, $C_L \to C_0/D$~: la première goutte
de liquide peut enrichir un élément au plus d'un facteur $1/D$. Pour $D \to 0$, $C_L \to
C_0/F$~: un élément parfaitement incompatible est dilué dans le liquide, et son
enrichissement ne mesure que $F$. Pour un élément compatible, $D > 1$, le liquide est
plus pauvre que la source, et d'autant plus que $F$ est petit.

\begin{proposition}[Fusion à l'équilibre non modale]\label{prop:b3:trace-element-modelling:nonmodal}
Dans la \omterm{thm:b3:trace-element-modelling:batch}{fusion à l'équilibre}, avec les proportions de fusion $p_i$ de la \cref{def:b3:trace-element-modelling:nonmodal},
\[
\frac{C_L}{C_0} = \frac{1}{D_0 + F(1 - P)}.
\]
\end{proposition}

\begin{proof}
D'après la \cref{prop:b3:trace-element-modelling:residue}, le résidu contient $\bar D(F)\,C_L$,
donc $C_0 = FC_L + (1 - F)\bar D(F)C_L = C_L(F + D_0 - PF)$.
\end{proof}

Lorsque $P$ dépasse $D_0$, comme pour le lanthane dans l'exemple traité, où $D_0 =
0.0074$ et $P = 0.025$, le résidu perd son meilleur hôte, le clinopyroxène, plus
vite que les autres. Le liquide est alors un peu plus riche que ne le prédit la fusion modale,
d'un peu moins de \qty{2}{\percent} à $F = 0.1$ pour le lanthane (exercice~6), mais le
résidu épuise son clinopyroxène à $F = 0.12/0.45 \approx 0.27$, et au-delà
de ce seuil, les minéraux et les coefficients changent.
% ledger: white13:ex73

\section{Fusion fractionnée et fusion dynamique}\label{sec:b3:trace-element-modelling:fractional}

Le liquide de fusion d'une péridotite se rassemble le long des joints de grains et peut
s'échapper dès qu'il est connecté. Le pôle opposé à la \omterm{thm:b3:trace-element-modelling:batch}{fusion à l'équilibre} retire chaque goutte
au fur et à mesure qu'elle se forme.

\begin{theorem}[Fusion fractionnée]\label{thm:b3:trace-element-modelling:fractional}
Dans la \emph{fusion fractionnée}\index{fusion fractionnée}, chaque incrément de liquide
est en équilibre avec le résidu seulement au moment où il se forme, puis le
quitte. Pour un $D$ constant, après extraction d'une masse totale $F$ de liquide, le
résidu et le liquide qui se forme à cet instant vérifient
\[
\frac{C_S}{C_0} = (1 - F)^{1/D - 1}, \qquad
\frac{C_L}{C_0} = \frac{1}{D}(1 - F)^{1/D - 1}.
\]
\end{theorem}

\begin{proof}
Le résidu a une masse $1 - F$ et une concentration $C_S$. Lorsqu'une masse
$\mathrm dF$ de liquide part à $C_L = C_S/D$, l'élément du résidu
diminue de $\mathrm d[(1 - F)C_S] = -(C_S/D)\,\mathrm dF$. En développant,
$(1 - F)\,\mathrm dC_S = C_S(1 - 1/D)\,\mathrm dF$, donc
$\mathrm dC_S/C_S = (1/D - 1)\,\mathrm d[-\ln(1 - F)]$. En intégrant depuis
$C_S = C_0$ en $F = 0$ donne $\ln(C_S/C_0) = (1/D - 1)\ln(1 - F)$, et
$C_L = C_S/D$.
\end{proof}

Pour un élément incompatible, $1/D - 1$ est grand et positif~: le résidu est
épuisé après quelques multiples de $D$ de fusion, et le liquide instantané passe
de $C_0/D$ à presque rien. Un résidu de \omterm{thm:b3:trace-element-modelling:fractional}{fusion fractionnée} est bien plus
appauvri qu'un résidu de \omterm{thm:b3:trace-element-modelling:batch}{fusion à l'équilibre}. Les liquides ne sont pas émis un à un,
toutefois~: ils se rassemblent dans des chenaux et forment des poches, et le liquide regroupé, défini ci-dessous,
est tracé en tirets sur la figure.

\begin{omfigure}
\begin{tikzpicture}
\begin{groupplot}[group style={group size=2 by 1, horizontal sep=1.6cm},
  omrecord, width=0.47\linewidth, height=0.44\linewidth, title style={font=\small},
  xlabel={taux de fusion $F$}, xmin=0, xmax=0.3, ymode=log, ymin=0.01, ymax=150,
  xtick={0,0.1,0.2,0.3}, xticklabels={0,0.1,0.2,0.3},
  ytick={0.01,0.1,1,10,100}, yticklabels={0.01,0.1,1,10,100},
  legend style={font=\scriptsize, legend cell align=left}]
\nextgroupplot[title={(a) fusion à l'équilibre (plein), fusion fractionnée agrégée (tirets)},
  ylabel={$C_L/C_0$}, legend style={at={(0.97,0.97)}, anchor=north east}]
% computed in figdata/bachelor-3/melting_models.py
\addplot[omDef, thick] table[x=f, y=b001] {figdata/out/bachelor-3-melting_models-curves.dat};
\addlegendentry{$D = 0.01$}
\addplot[red!70!black, thick] table[x=f, y=b01] {figdata/out/bachelor-3-melting_models-curves.dat};
\addlegendentry{$D = 0.1$}
\addplot[black, thick] table[x=f, y=b1] {figdata/out/bachelor-3-melting_models-curves.dat};
\addlegendentry{$D = 1$}
\addplot[omLithLime!70!black, thick] table[x=f, y=b10] {figdata/out/bachelor-3-melting_models-curves.dat};
\addlegendentry{$D = 10$}
\addplot[omDef, thick, densely dashed] table[x=f, y=a001] {figdata/out/bachelor-3-melting_models-curves.dat};
\addplot[red!70!black, thick, densely dashed] table[x=f, y=a01] {figdata/out/bachelor-3-melting_models-curves.dat};
\addplot[omLithLime!70!black, thick, densely dashed] table[x=f, y=a10] {figdata/out/bachelor-3-melting_models-curves.dat};
\nextgroupplot[title={(b) fractionnée, instantanée}]
\addplot[omDef, thick] table[x=f, y=i001] {figdata/out/bachelor-3-melting_models-curves.dat};
\addplot[red!70!black, thick] table[x=f, y=i01] {figdata/out/bachelor-3-melting_models-curves.dat};
\addplot[black, thick] table[x=f, y=i1] {figdata/out/bachelor-3-melting_models-curves.dat};
\addplot[omLithLime!70!black, thick] table[x=f, y=i10] {figdata/out/bachelor-3-melting_models-curves.dat};
\end{groupplot}
\end{tikzpicture}

\par\smallskip{\small Enrichissement d'un élément en traces dans le liquide en fonction du
taux de fusion, pour quatre coefficients de partage globaux. (a) Liquides de \omterm{thm:b3:trace-element-modelling:batch}{fusion à l'équilibre} (plein)
et \omterm{thm:b3:trace-element-modelling:fractional}{fusion fractionnée} agrégée (tirets~; pour $D = 1$ les deux restent à 1). (b) Les
liquides instantanés de la \omterm{thm:b3:trace-element-modelling:fractional}{fusion fractionnée}~: un élément avec $D = 0.01$ est
épuisé du résidu après quelques pour cent de fusion.}\label{fig:b3:trace-element-modelling:curves}
\end{omfigure}

\begin{definition}[Liquide de fusion fractionnée agrégé]\label{def:b3:trace-element-modelling:fractional}
Le \emph{liquide de fusion fractionnée agrégé}\index{liquide de fusion fractionnée agrégé} est le
mélange de tous les liquides instantanés d'une \omterm{thm:b3:trace-element-modelling:fractional}{fusion fractionnée}, de $0$ à
$F$, regroupés sans réaction ultérieure avec le résidu.
\end{definition}

\begin{proposition}[Composition du liquide agrégé]
\label{prop:b3:trace-element-modelling:aggregate}
Pour un $D$ constant, le \omterm{def:b3:trace-element-modelling:fractional}{liquide de fusion fractionnée agrégé} vérifie
\[
\frac{\bar C_L}{C_0} = \frac{1 - (1 - F)^{1/D}}{F}.
\]
\end{proposition}

\begin{proof}
L'élément extrait est la source moins le résidu,
$C_0 - (1 - F)C_S = C_0[1 - (1 - F)^{1/D}]$ d'après le \cref{thm:b3:trace-element-modelling:fractional},
réparti sur une masse $F$ de liquide.
\end{proof}

Dès que $(1 - F)^{1/D}$ est petit, le liquide agrégé contient la totalité de
l'élément et $\bar C_L \approx C_0/F$, quel que soit $D$~: pour les éléments très
incompatibles, les \omterm{def:b3:trace-element-modelling:fractional}{liquides de fusion fractionnée agrégés} et les liquides de \omterm{thm:b3:trace-element-modelling:batch}{fusion
à l'équilibre} à petit $D$ coïncident, et mesurent tous deux $F$.

La fusion réelle se situe entre les deux pôles. Un résidu conserve une partie du
liquide dans ses pores, et le liquide qui quitte la source est l'excès au-delà de cette porosité.

\begin{definition}[Fusion dynamique]\label{def:b3:trace-element-modelling:dynamic}
Dans la \emph{fusion dynamique}\index{fusion dynamique} (appelée aussi fusion
continue), le résidu conserve une fraction massique constante $\phi$ de liquide
dans ses pores, en équilibre avec ses cristaux, et tout liquide produit au-delà
est extrait.
\end{definition}

\begin{proposition}[Fusion dynamique]\label{prop:b3:trace-element-modelling:dynamic}
La \omterm{def:b3:trace-element-modelling:dynamic}{fusion dynamique} de porosité $\phi$ est une \omterm{thm:b3:trace-element-modelling:fractional}{fusion fractionnée} de coefficient de
partage effectif $D' = (1 - \phi)D + \phi$, où $F$ est la fraction massique de
liquide extrait~:
\[
\frac{C_L}{C_0} = \frac{1}{D'}(1 - F)^{1/D' - 1}.
\]
\end{proposition}

\begin{proof}
Le résidu est constitué de cristaux, de fraction $1 - \phi$, et de liquide de pores,
de fraction $\phi$, tous deux en équilibre avec le liquide de concentration $C_L$, donc
il vérifie $C_S = [(1 - \phi)D + \phi]C_L = D'C_L$. La démonstration du
\cref{thm:b3:trace-element-modelling:fractional} s'applique alors avec $D'$.
\end{proof}

Pour $\phi = 0$, le modèle est celui de la \omterm{thm:b3:trace-element-modelling:fractional}{fusion fractionnée}~; pour $\phi \to 1$, il
tend vers la \omterm{thm:b3:trace-element-modelling:batch}{fusion à l'équilibre}. Une porosité de un pour cent plafonne
l'enrichissement du premier liquide à $1/\phi = 100$, quel que soit $D$~: le liquide
de pores tamponne le résidu.

\begin{omfigure}
\begin{tikzpicture}[font=\small, >=Stealth]
\tikzset{rock/.style={draw, fill=omLithLime!25, minimum width=2.3cm, minimum height=2.0cm},
  melt/.style={draw, fill=red!25}}
% (a) batch
\node[rock] (a) at (0,0) {};
\foreach \x/\y in {-0.7/0.5, 0.3/0.6, -0.2/-0.1, 0.7/-0.4, -0.6/-0.6, 0.4/0.0}
  \fill[red!55] (\x,\y) circle (0.09);
\node[align=center, font=\scriptsize] at (0,-1.45) {résidu $1-F$ à $DC_L$\\liquide $F$ à $C_L$};
\draw[->, thick, red!70!black] (0,1.05) -- (0,1.7) node[above, font=\scriptsize] {d'un coup};
\node[font=\small] at (0,-2.25) {(a) à l'équilibre};
% (b) fractional
\node[rock] (b) at (4,0) {};
\draw[->, thick, red!70!black] (4,1.05) -- (4,1.7);
\foreach \y/\s in {1.85/0.06, 2.05/0.08, 2.27/0.1}
  \fill[red!55] (4,\y) circle (\s);
\node[align=center, font=\scriptsize] at (4,-1.45) {chaque goutte part\\dès qu'elle se forme};
\node[font=\small] at (4,-2.25) {(b) fractionnée};
% (c) dynamic
\node[rock] (c) at (8,0) {};
\foreach \x/\y in {7.4/0.5, 8.4/0.3, 7.8/-0.5, 8.6/-0.6}
  \fill[red!55] (\x,\y) circle (0.06);
\draw[->, thick, red!70!black] (8,1.05) -- (8,1.7);
\foreach \y/\s in {1.85/0.06, 2.05/0.08, 2.27/0.1}
  \fill[red!55] (8,\y) circle (\s);
\node[align=center, font=\scriptsize] at (8,-1.45) {les pores gardent $\phi$~;\\l'excès part};
\node[font=\small] at (8,-2.25) {(c) dynamique};
\end{tikzpicture}

\par\smallskip{\small Trois façons pour un liquide de quitter une source en fusion. (a)
À l'équilibre~: le liquide reste avec tout le résidu, puis part. (b) Fractionnée~:
chaque incrément part aussitôt~; les incréments regroupés forment le liquide
agrégé. (c) Dynamique~: le résidu conserve une petite fraction massique $\phi$ de
liquide dans ses pores.}\label{fig:b3:trace-element-modelling:modes}
\end{omfigure}

\section{Fusion, puis cristallisation}\label{sec:b3:trace-element-modelling:crystallisation}

Un liquide qui quitte sa source ne sort que rarement inchangé. Il cristallise
de l'olivine, puis du pyroxène et du plagioclase sur son trajet vers la surface~;
il peut assimiler la roche encaissante, et se mélanger à d'autres lots. Les lois du volume de deuxième année s'appliquent
tour à tour~: une lave est un liquide de fraction $F$ de sa source, puis un liquide qui
a conservé une fraction $f$ de sa masse par cristallisation fractionnée, de sorte que
pour un élément de coefficient $D$ dans le résidu de fusion et $D_c$ dans les
minéraux qui cristallisent,
\[
C_{\text{lave}} = C_0\,\frac{1}{D + F(1 - D)}\,f^{\,D_c - 1}.
\]
Le second facteur est le même pour tous les éléments que les minéraux qui cristallisent
rejettent, $D_c \approx 0$~: il les multiplie tous par $1/f$. Les rapports entre de tels
éléments sont donc fixés à la source, et une lave ayant perdu de l'olivine conserve le
rapport La/Yb de son liquide primaire tandis que son lanthane augmente. Inverser $F$
à partir d'une concentration exige aussi $f$~; inverser la source à partir d'un rapport
ne l'exige pas. Le plagioclase fait exception parmi les terres rares~: il capte
l'europium, et son retrait laisse l'anomalie en europium du volume de deuxième année.

\section{Lire les profils de terres rares}\label{sec:b3:trace-element-modelling:patterns}

Un liquide de lherzolite à spinelle et un liquide de lherzolite à grenat partagent
leurs terres rares légères, que ni l'un ni l'autre résidu ne retient, et diffèrent
par les lourdes, que le grenat retient. Sur un diagramme normalisé aux chondrites,
les liquides de grenat sont raides et à palier dans la partie lourde~; les liquides de spinelle sont plus doux. La
raideur croît à mesure que $F$ diminue.

\begin{omfigure}
\begin{tikzpicture}
\begin{axis}[omrecord, width=0.8\linewidth, height=0.5\linewidth, ymode=log,
  xmin=56.5, xmax=71.5, ymin=1, ymax=1000, xtick={57,58,60,62,63,64,65,70,71},
  xticklabels={La,Ce,Nd,Sm,Eu,Gd,Tb,Yb,Lu}, ytick={1,10,100,1000}, yticklabels={1,10,100,1000},
  ylabel={échantillon / chondrite CI},
  legend style={font=\scriptsize, at={(0.98,0.98)}, anchor=north east,
  legend cell align=left, legend columns=2}]
% computed in figdata/bachelor-3/melting_models.py
\addplot[omDef, thick, mark=*, mark size=1.3pt] table[x=z, y=sp01] {figdata/out/bachelor-3-melting_models-ree.dat};
\addlegendentry{spinelle, $F = 1\,\%$}
\addplot[red!70!black, thick, mark=square*, mark size=1.3pt] table[x=z, y=gt01] {figdata/out/bachelor-3-melting_models-ree.dat};
\addlegendentry{grenat, $F = 1\,\%$}
\addplot[omDef, thick, densely dashed, mark=*, mark size=1.3pt] table[x=z, y=sp05] {figdata/out/bachelor-3-melting_models-ree.dat};
\addlegendentry{spinelle, $5\,\%$}
\addplot[red!70!black, thick, densely dashed, mark=square*, mark size=1.3pt] table[x=z, y=gt05] {figdata/out/bachelor-3-melting_models-ree.dat};
\addlegendentry{grenat, $5\,\%$}
\addplot[omDef, thick, densely dotted, mark=*, mark size=1.3pt] table[x=z, y=sp20] {figdata/out/bachelor-3-melting_models-ree.dat};
\addlegendentry{spinelle, $20\,\%$}
\addplot[red!70!black, thick, densely dotted, mark=square*, mark size=1.3pt] table[x=z, y=gt20] {figdata/out/bachelor-3-melting_models-ree.dat};
\addlegendentry{grenat, $20\,\%$}
\addplot[black, thin, mark=none] table[x=z, y=pm] {figdata/out/bachelor-3-melting_models-ree.dat};
\addlegendentry{source}
\end{axis}
\end{tikzpicture}

\par\smallskip{\small Profils de terres rares normalisés aux chondrites de liquides de \omterm{thm:b3:trace-element-modelling:batch}{fusion
à l'équilibre} d'une source de manteau primitif, issus d'une lherzolite à spinelle (cercles bleus) et d'une
lherzolite à grenat (carrés rouges), à trois taux de fusion~; la source est la
ligne fine. Seules les neuf terres rares dont les coefficients de partage sont tabulés sont
représentées.}\label{fig:b3:trace-element-modelling:ree}
\end{omfigure}
% ledger: budahn91:d-ree, ms95:pm-trace, lodders:ci-table, white13:ex73, rfr25:gkr-5gpa

Le grenat cède la place au spinelle comme minéral alumineux du manteau vers 60
à \qty{90}{km} de profondeur. Un liquide dont les terres rares lourdes sont plates et basses,
avec une pente forte des légères vers les lourdes, a donc laissé du grenat dans son résidu et
a commencé à fondre à une profondeur d'au moins celle-ci.
% ledger: white13:sp-gt

Un profil mélange deux inconnues, $F$ et la source. Porter deux éléments
l'un contre l'autre permet de les séparer.

\begin{definition}[Diagramme rapport--rapport]\label{def:b3:trace-element-modelling:ratios}
Un \emph{diagramme rapport--rapport}\index{diagramme rapport--rapport} porte, pour une suite
de laves, le rapport d'un élément très incompatible $H$ à un élément moins
incompatible $M$ en fonction de la concentration en $H$ (ou en fonction d'un autre
rapport). La fusion partielle, la cristallisation et le mélange tracent sur lui des courbes
de formes différentes.
\end{definition}

\begin{proposition}[La droite de fusion à l'équilibre]\label{prop:b3:trace-element-modelling:line}
Pour les liquides de \omterm{thm:b3:trace-element-modelling:batch}{fusion à l'équilibre} d'une même source avec $D_H = 0$ et un $D_M$ constant, les points
$(C_H, C_H/C_M)$ se placent sur la droite
\[
\frac{C_H}{C_M} = \frac{D_M}{C_0^M}\,C_H + \frac{C_0^H(1 - D_M)}{C_0^M}.
\]
\end{proposition}

\begin{proof}
Avec $D_H = 0$, $C_H = C_0^H/F$, donc $1/F = C_H/C_0^H$. Alors
$C_H/C_M = (C_0^H/F)\,(D_M + F(1 - D_M))/C_0^M = C_0^HD_M/(FC_0^M) +
C_0^H(1 - D_M)/C_0^M$, et le premier terme vaut $(D_M/C_0^M)\,C_H$.
\end{proof}

La cristallisation fractionnée de minéraux qui rejettent les deux éléments déplace une lave
horizontalement sur ce diagramme, puisque le rapport ne change pas~; la fusion la déplace
le long d'une droite inclinée dont la pente croît avec $D_M$~; le mélange de deux liquides dessine une
hyperbole. Une source à grenat a un $D_{\mathrm{Yb}}$ élevé, donc une droite de fusion
raide dans le diagramme La/Yb en fonction de La.

\begin{method}[Inversion de deux laves]\label{met:b3:trace-element-modelling:invert}
(1) Vérifier que les laves sont cogénétiques et que $H$ est bien plus incompatible
que $M$. (2) Corriger chaque lave de la cristallisation, ou n'utiliser que des rapports.
(3) À partir de deux laves ou plus, ajuster la droite $C_H/C_M = aC_H + b$.
(4) Alors $a = D_M/C_0^M$ et $b = C_0^H(1 - D_M)/C_0^M$~: $C_0^M$ étant supposé, $D_M =
aC_0^M$ et $C_0^H = bC_0^M/(1 - D_M)$. (5) Le taux de fusion de chaque lave est
$F = C_0^H/C_H$.
\end{method}

\begin{omfigure}
\begin{tikzpicture}
\begin{axis}[omrecord, width=0.78\linewidth, height=0.52\linewidth,
  xmode=log, ymode=log, xmin=2, xmax=60, ymin=2, ymax=100,
  xlabel={La (\unit{ppm})}, ylabel={La/Yb},
  xtick={2,5,10,20,50}, xticklabels={2,5,10,20,50},
  ytick={2,5,10,20,50,100}, yticklabels={2,5,10,20,50,100},
  legend style={font=\scriptsize, at={(0.03,0.97)}, anchor=north west,
  legend cell align=left}]
% computed in figdata/bachelor-3/melting_models.py
\addplot[omDef, thick] table[x=la_sp, y=lay_sp] {figdata/out/bachelor-3-melting_models-traj.dat};
\addlegendentry{lherzolite à spinelle}
\addplot[red!70!black, thick] table[x=la_gt, y=lay_gt] {figdata/out/bachelor-3-melting_models-traj.dat};
\addlegendentry{lherzolite à grenat}
\addplot[omDef, only marks, mark=|, mark size=3pt, forget plot] table[x=la_sp, y=lay_sp] {figdata/out/bachelor-3-melting_models-ticks.dat};
\addplot[red!70!black, only marks, mark=|, mark size=3pt, forget plot] table[x=la_gt, y=lay_gt] {figdata/out/bachelor-3-melting_models-ticks.dat};
\addplot[black, only marks, mark=diamond*, mark size=2.8pt] table[x=la, y=layb] {figdata/out/bachelor-3-melting_models-lavas.dat};
\addlegendentry{laves A et B}
\node[font=\scriptsize, anchor=north] at (axis cs:9.6,5.2) {A};
\node[font=\scriptsize, anchor=south east] at (axis cs:15.2,36) {B};
\node[font=\scriptsize, omDef, anchor=north west] at (axis cs:36.5,17.5) {1\,\%};
\node[font=\scriptsize, omDef, anchor=north west] at (axis cs:11.3,6.2) {5\,\%};
\node[font=\scriptsize, omDef, anchor=north west] at (axis cs:3.14,2.55) {20\,\%};
\node[font=\scriptsize, red!70!black, anchor=south east] at (axis cs:29.5,66) {1\,\%};
\node[font=\scriptsize, red!70!black, anchor=south east] at (axis cs:10.5,24) {5\,\%};
\node[font=\scriptsize, red!70!black, anchor=south east] at (axis cs:3.09,7.0) {20\,\%};
\end{axis}
\end{tikzpicture}

\par\smallskip{\small Liquides de \omterm{thm:b3:trace-element-modelling:batch}{fusion à l'équilibre} d'une source de manteau primitif dans un diagramme La/Yb
en fonction de La, pour une lherzolite à spinelle et une lherzolite à grenat, avec des repères à $F = 1$,
2, 5, 10 et \qty{20}{\percent}, et les deux laves fictives du
problème du week-end. Chaque lave ne se situe que sur une seule courbe.}\label{fig:b3:trace-element-modelling:ratios}
\end{omfigure}

En échelle logarithmique, la droite de la proposition s'incurve, et avec
$D_{\mathrm{La}} \ne 0$ elle n'est presque droite qu'en échelle linéaire~; les courbes
de la figure sont exactes. Une source à spinelle ne peut pas produire un liquide de rapport La/Yb supérieur à
$(C_0^{\mathrm{La}}/D_{\mathrm{La}})/(C_0^{\mathrm{Yb}}/D_{\mathrm{Yb}})$,
soit environ 40 pour le manteau primitif~; une source à grenat peut atteindre environ 120.

\section{Exercices}

\begin{exercise}[$\star$]\label{exo:b3:trace-element-modelling:1}
Avec les coefficients minéraux du samarium, 0.40 dans le clinopyroxène, 0.11
dans le grenat et 0.016 dans l'olivine, l'orthopyroxène et le spinelle, calculer le
coefficient de partage global du samarium dans la lherzolite à spinelle (\qty{12}{\percent}
de clinopyroxène) et dans la lherzolite à grenat (\qty{19.5}{\percent} de clinopyroxène,
\qty{16.7}{\percent} de grenat) du chapitre.
\end{exercise}
% ledger: budahn91:d-ree

\begin{exercise}[$\star$]\label{exo:b3:trace-element-modelling:2}
Un élément de $D = 0.05$ est présent à \qty{1}{ppm} dans une source. Calculer
sa concentration dans les liquides de \omterm{thm:b3:trace-element-modelling:batch}{fusion à l'équilibre} pour $F = 0.02$ et $0.10$.
\end{exercise}

\begin{exercise}[$\star$]\label{exo:b3:trace-element-modelling:3}
Quel est le plus fort enrichissement $C_L/C_0$ que la \omterm{thm:b3:trace-element-modelling:batch}{fusion à l'équilibre} peut
donner à un élément de $D = 0.01$~? Et à un élément de $D = 0$ pour
$F = \qty{5}{\percent}$~?
\end{exercise}

\begin{exercise}[$\star$]\label{exo:b3:trace-element-modelling:4}
Une lave contient \qty{30}{ppm} de La et \qty{2.5}{ppm} de Yb. Avec les valeurs
de la chondrite CI, La 0.232 et Yb \qty{0.163}{ppm}, calculer son rapport La/Yb et son
rapport normalisé à la chondrite $(\mathrm{La}/\mathrm{Yb})_N$.
\end{exercise}
% ledger: lodders:ci-table

\begin{exercise}[$\star\star$]\label{exo:b3:trace-element-modelling:5}
Pour un élément de $D = 0.01$, comparer les résidus après \qty{5}{\percent}
de \omterm{thm:b3:trace-element-modelling:batch}{fusion à l'équilibre} et de \omterm{thm:b3:trace-element-modelling:fractional}{fusion fractionnée}. Que contient le liquide de \omterm{thm:b3:trace-element-modelling:fractional}{fusion
fractionnée} instantané à cet instant~?
\end{exercise}

\begin{exercise}[$\star\star$]\label{exo:b3:trace-element-modelling:6}
Dans l'exemple traité du chapitre (proportions modales 58, 27, 12, \qty{3}{\percent}~;
proportions de fusion 20, 25, 45, \qty{10}{\percent} pour l'olivine,
l'orthopyroxène, le clinopyroxène et le spinelle), le lanthane a pour coefficients
0.0003, 0.002, 0.053 et 0.01. Calculer $D_0$ et $P$, puis l'enrichissement du
lanthane dans le liquide de \omterm{thm:b3:trace-element-modelling:batch}{fusion à l'équilibre} pour $F = 0.02$ et $0.10$, en \omterm{def:b3:trace-element-modelling:nonmodal}{fusion non modale} puis modale.
\end{exercise}
% ledger: white13:ex73

\begin{exercise}[$\star\star$]\label{exo:b3:trace-element-modelling:7}
Pour $F = 0.10$, comparer le liquide de \omterm{thm:b3:trace-element-modelling:batch}{fusion à l'équilibre} et le \omterm{def:b3:trace-element-modelling:fractional}{liquide de fusion
fractionnée agrégé} pour $D = 0.1$ puis pour $D = 0.01$. Lequel diffère le plus, et pourquoi~?
\end{exercise}

\begin{exercise}[$\star\star$]\label{exo:b3:trace-element-modelling:8}
Un élément de $D = 0.001$ fond dynamiquement avec $\phi = 0.01$. Calculer
$D'$, l'enrichissement du premier liquide extrait, et celui du liquide extrait
à $F = 0.01$~; comparer avec la \omterm{thm:b3:trace-element-modelling:fractional}{fusion fractionnée}.
\end{exercise}

\begin{exercise}[$\star\star$]\label{exo:b3:trace-element-modelling:9}
Dans la lherzolite à spinelle du chapitre, $D_{\mathrm{La}} = 0.0078$ et
$D_{\mathrm{Sm}} = 0.062$. Par quel facteur la \omterm{thm:b3:trace-element-modelling:batch}{fusion à l'équilibre} augmente-t-elle le rapport La/Sm
au-dessus de la valeur de la source pour $F = 0.01$ et $F = 0.10$~? Qu'est-ce que cela
indique sur l'usage des rapports pour estimer $F$~?
\end{exercise}

\begin{exercise}[$\star\star\star$]\label{exo:b3:trace-element-modelling:10}
Montrer que le \omterm{def:b3:trace-element-modelling:fractional}{liquide de fusion fractionnée agrégé} est la moyenne des liquides
instantanés de 0 à $F$, et qu'il tend vers $C_0/F$ lorsque $(1 - F)^{1/D} \to 0$.
Pour $D = 0.01$, à partir de quel $F$ est-il à moins de \qty{1}{\percent} de $C_0/F$~?
\end{exercise}

\begin{exercise}[$\star\star\star$]\label{exo:b3:trace-element-modelling:11}
Deux laves cogénétiques P et Q ont un La de 20 et \qty{10}{ppm}, et des rapports La/Yb de 10 et 6.
En supposant le lanthane parfaitement incompatible et une source contenant
\qty{0.44}{ppm} de Yb, appliquer la méthode du chapitre~: trouver $D_{\mathrm{Yb}}$,
le La de la source et le taux de fusion de chaque lave.
\end{exercise}

\begin{exercise}[$\star\star\star$]\label{exo:b3:trace-element-modelling:12}
Dans la \omterm{def:b3:trace-element-modelling:nonmodal}{fusion non modale} de l'exercice~6, à quel taux de fusion chaque minéral
est-il épuisé~? Calculer le coefficient global du résidu pour le lanthane à
$F = 0.2$, et dire pourquoi la formule du chapitre ne peut pas servir au-delà
du premier épuisement.
\end{exercise}

\section{Problème~: deux laves de l'île d'Aldren}

\begin{problem}[{Problème du week-end --- la source, la profondeur et le taux de fusion
de deux basaltes fictifs, à partir de leur lanthane et de leur ytterbium}]\label{pb:b3:trace-element-modelling:1}
Sur l'île d'Aldren, île volcanique fictive, la lave A contient \qty{9.6}{ppm} de La
et \qty{1.7}{ppm} de Yb, la lave B \qty{15.6}{ppm} de La et \qty{0.45}{ppm} de Yb.
Toutes deux sont prises comme liquides primaires d'une source de manteau primitif
contenant \qty{0.648}{ppm} de La et \qty{0.441}{ppm} de Yb (chondrite CI~: La 0.232, Yb
\qty{0.163}{ppm}). Coefficients minéraux~: La 0.058 dans le clinopyroxène et 0.001
dans le grenat, l'olivine, l'orthopyroxène et le spinelle~; Yb 0.58 dans le
clinopyroxène, 4.5 dans le grenat et 0.16 dans l'olivine, l'orthopyroxène et le spinelle.
La lherzolite à spinelle comprend \qty{58}{\percent} d'olivine, \qty{27}{\percent} d'orthopyroxène, \qty{12}{\percent}
de clinopyroxène, \qty{3}{\percent} de spinelle~; la lherzolite à grenat \qty{54.1}{\percent} d'olivine, \qty{9.8}{\percent}
d'orthopyroxène, \qty{19.5}{\percent} de clinopyroxène, \qty{16.7}{\percent} de grenat. La fusion est modale.
% ledger: ms95:pm-trace, lodders:ci-table, budahn91:d-ree, white13:ex73, rfr25:gkr-5gpa

\medskip\noindent\textbf{Partie I --- Les deux sources.}
\begin{enumerate}
\item Calculer $D_{\mathrm{La}}$ et $D_{\mathrm{Yb}}$ de la lherzolite à spinelle.
\item Les calculer pour la lherzolite à grenat. Arrondir les quatre valeurs à deux chiffres
significatifs pour la suite du problème.
\item Quel coefficient diffère le plus entre les sources, et pourquoi~? L'ytterbium est-il
compatible dans la lherzolite à grenat~?
\item Calculer le $(\mathrm{La}/\mathrm{Yb})_N$ normalisé à la chondrite de la
source. À quoi ressemble son spectre~?
\end{enumerate}

\medskip\noindent\textbf{Partie II --- Fusion à l'équilibre de la lave A.}
\begin{enumerate}[resume]
\item Calculer La dans les liquides de \omterm{thm:b3:trace-element-modelling:batch}{fusion à l'équilibre} de la lherzolite à spinelle pour $F = 1$, 5 et
\qty{10}{\percent}.
\item Trouver le taux de fusion qui donne à la lave A son lanthane.
\item Prédire son ytterbium et son La/Yb, et comparer avec la lave A.
\item Calculer La et Yb dans le résidu, ainsi que la fraction du lanthane de la source
lanthane qui y reste.
\item Quel taux de fusion trouverait-on en prenant le lanthane comme parfaitement
incompatible~?
\end{enumerate}

\medskip\noindent\textbf{Partie III --- Fusion fractionnée de la lave A.}
\begin{enumerate}[resume]
\item Calculer La dans le liquide de \omterm{thm:b3:trace-element-modelling:fractional}{fusion fractionnée} instantané pour $F = 0.1$, 1 et
\qty{6}{\percent}.
\item Calculer La dans le \omterm{def:b3:trace-element-modelling:fractional}{liquide de fusion fractionnée agrégé} pour $F = \qty{6.0}{\percent}$.
\item Trouver le taux de fusion auquel le liquide agrégé possède le
lanthane de la lave A. Pourquoi est-il si proche de la réponse à la question~9~?
\item Comparer le lanthane du résidu avec celui de la question~8.
\item Quelle estimation de $F$ retenir~?
\end{enumerate}

\medskip\noindent\textbf{Partie IV --- La lave B et le grenat.}
\begin{enumerate}[resume]
\item Calculer La/Yb et $(\mathrm{La}/\mathrm{Yb})_N$ des deux laves.
\item Quel est le La/Yb le plus élevé qu'un liquide de \omterm{thm:b3:trace-element-modelling:batch}{fusion à l'équilibre} de la lherzolite à spinelle
peut atteindre~?
\item À quel taux de fusion un liquide de lherzolite à spinelle aurait-il le La/Yb de
la lave B, et combien de lanthane contiendrait-il~? Conclure.
\item Trouver le taux de fusion de la lave B à partir de son lanthane, pour la lherzolite à grenat,
et prédire son ytterbium.
\item Que disent les deux sources des profondeurs auxquelles A et B ont commencé à
fondre~? La lave A pourrait-elle provenir de la lherzolite à grenat~?
\item Calculer la pente et l'ordonnée à l'origine de la droite de \omterm{thm:b3:trace-element-modelling:batch}{fusion à l'équilibre} La/Yb $= a\,
\mathrm{La} + b$ de la source à grenat, en prenant $D_{\mathrm{La}}$ nul,
et vérifier la lave B par rapport à celle-ci.
\item Si la lave A avait perdu \qty{10}{\percent} de sa masse en olivine avant
son éruption, comment changeraient son La/Yb et son taux de fusion~?
\item Répondre à la question d'ouverture du chapitre pour ces deux laves.
\item Donner le degré de fusion de la lave A.
\end{enumerate}
\end{problem}
